\documentclass[12pt, a4paper]{article}

% ─── Packages ────────────────────────────────────────────────────────────────
\usepackage[T1]{fontenc}
\usepackage[utf8]{inputenc}
\usepackage[english]{babel}
\usepackage{geometry}
\usepackage{graphicx}
\usepackage{amsmath}
\usepackage{amssymb}
\usepackage{siunitx}
\usepackage{booktabs}
\usepackage{caption}
\usepackage{subcaption}
\usepackage{float}
\usepackage{hyperref}
\usepackage{fancyhdr}
\usepackage{parskip}        % automatic paragraph spacing, no manual \\
\usepackage{xcolor}
\usepackage{array}
\usepackage{multirow}
\usepackage{enumitem}
\usepackage{tcolorbox}
\usepackage{tabularx}
\usepackage{microtype}      % better line breaking

\geometry{top=2.5cm, bottom=2.5cm, left=2.5cm, right=2.5cm}

% ─── siunitx ─────────────────────────────────────────────────────────────────
\sisetup{per-mode=symbol, detect-weight=true}

% ─── Placeholder commands ────────────────────────────────────────────────────
% Measured value placeholder:  \val{unit}  ->  blue bold "?? unit"
\newcommand{\val}[1]{\textcolor{blue}{\textbf{??\,\si{#1}}}}

% Measured value without unit:  \valx  ->  blue bold "??"
\newcommand{\valx}{\textcolor{blue}{\textbf{??}}}

% Free-text placeholder:  \enter{description}  ->  blue italic [description]
\newcommand{\enter}[1]{\textcolor{blue}{\textit{[#1]}}}

% Hint box:  \hint{text}  ->  grey box with bullet hints
\tcbuselibrary{skins}
\newcommand{\hint}[1]{%
  \begin{tcolorbox}[
    enhanced,
    colback  = gray!7,
    colframe = gray!35,
    fontupper= \small,
    left=6pt, right=6pt, top=4pt, bottom=4pt,
    title    = \textbf{\small Hints for this section},
    attach boxed title to top left={yshift=-2mm, xshift=4mm},
    boxed title style={colback=gray!35, colframe=gray!35,
                       fontupper=\small\bfseries, left=3pt, right=3pt}
  ]
  #1
  \end{tcolorbox}}

% ─── Hyperref ────────────────────────────────────────────────────────────────
\hypersetup{colorlinks=true, linkcolor=black, citecolor=black, urlcolor=blue,
            pdftitle={Lab Report - Switching Circuits},
            pdfauthor={Rupp Arian, Lampert Mathias}}

% ─── Header / Footer ─────────────────────────────────────────────────────────
\pagestyle{fancy}
\fancyhf{}
\lhead{\small EIT \& MEC -- Digi Lab 1: Switching Circuits}
\rhead{\small Fachhochschule Vorarlberg}
\cfoot{\small Page \thepage}
\lfoot{\small Rev.\,1.1 \quad 21/22.04.2026}
\renewcommand{\headrulewidth}{0.4pt}
\renewcommand{\footrulewidth}{0.4pt}

% ═════════════════════════════════════════════════════════════════════════════
\begin{document}
% ═════════════════════════════════════════════════════════════════════════════

% ─── Title page ──────────────────────────────────────────────────────────────
\begin{titlepage}
  \centering
  \vspace*{3cm}
  {\LARGE\bfseries Laboratory Report\par}
  \vspace{0.5cm}
  {\large\bfseries EIT \& MEC -- Digi Lab 1: Switching Circuits\par}
  \vspace{0.3cm}
  {\large Fachhochschule Vorarlberg GmbH\par}
  \vspace{2cm}
  \begin{tabular}{ll}
    \textbf{Students:}          & Rupp Arian, Lampert Mathias \\[4pt]
    \textbf{Study programme:}   & Elektronik und Informationstechnologie Dual \\[4pt]
    \textbf{Date of experiment:}& 21/22.04.2026 \\[4pt]
    \textbf{Supervisor:}        & Horatiu O.\,Pilsan \\[4pt]
    \textbf{Report guide:}      & C.\,Thorrold, \textit{Guide for Writing Lab Reports}, v1.0 \\
  \end{tabular}
  \vfill
  {\small Submitted: \today}
\end{titlepage}

% ─── Abstract ────────────────────────────────────────────────────────────────
\begin{abstract}
This report documents the experimental investigation of monostable and astable
multivibrators as well as a quartz crystal oscillator, conducted in Digi Lab~1 at
Fachhochschule Vorarlberg on 21/22~April~2026.
Four circuits were assembled and analysed: a monostable multivibrator based on the
NE555 timer IC, an astable multivibrator also using the NE555, an astable oscillator
built from two CMOS inverters (74C00), and a 12\,MHz quartz crystal oscillator.
For each circuit, theoretically calculated values for pulse duration and oscillation
period were compared against oscilloscope measurements, and deviations were discussed
with respect to component tolerances and measurement equipment effects.
\enter{Refine after experiments are completed -- max.\ 200 words.}
\end{abstract}

\newpage
\tableofcontents
\newpage

% ─────────────────────────────────────────────────────────────────────────────
\section{Introduction}
% ─────────────────────────────────────────────────────────────────────────────

Switching circuits are fundamental building blocks in digital and mixed-signal
electronics. This laboratory exercise focuses on two important classes of such
circuits: multivibrators, which generate or shape rectangular pulse signals, and
crystal oscillators, which provide highly stable reference frequencies.

A \textbf{monostable multivibrator} (one-shot) produces a single output pulse of
defined duration in response to a trigger event. An \textbf{astable multivibrator}
oscillates continuously between two states without any external trigger, generating
a periodic square wave. Both types can be realised with the NE555 timer IC, whose
internal structure consists of two comparators, a flip-flop, a discharge transistor,
and an output stage.

The \textbf{NE555} is a versatile and widely used timer IC. In monostable mode, the
output pulse width is determined solely by an external RC network according to
$t_H = 1.1\,R\,C$. In astable mode, the frequency and duty cycle depend on two
resistors and one capacitor.

As a complementary approach, an astable oscillator can also be constructed from two
cascaded CMOS inverter gates with a feedback RC network. Finally, a quartz crystal
oscillator exploits the piezoelectric resonance of a quartz crystal to achieve
frequency stabilities several orders of magnitude better than RC-based circuits.

The objectives of this laboratory are:
\begin{itemize}[leftmargin=1.5em]
  \item To understand and verify the operating principle of monostable and astable
        multivibrators using the NE555 IC.
  \item To compare calculated and measured values for pulse duration and oscillation
        period, and to quantify and explain deviations.
  \item To investigate the influence of the measurement setup (oscilloscope probes
        vs.\ BNC cables) on the circuit behaviour.
  \item To analyse the astable multivibrator realised with two CMOS inverters.
  \item To observe and explain the behaviour of a quartz crystal oscillator under
        varying supply voltages.
\end{itemize}

The theoretical background was reviewed according to~\cite{Floyd2006}
and~\cite{TietzeSchenk1999}.

% ─────────────────────────────────────────────────────────────────────────────
\section{Equipment and Materials}
% ─────────────────────────────────────────────────────────────────────────────

All circuits were assembled on a standard experimentation (breadboard) system.
Signals were observed and recorded with a digital storage oscilloscope.
Table~\ref{tab:equipment} lists the equipment used.

\begin{table}[H]
  \centering
  \caption{Equipment and components used in this experiment.}
  \label{tab:equipment}
  \begin{tabular}{llc}
    \toprule
    Item & Specification / Part No. & Asset / S/N \\
    \midrule
    Digital oscilloscope        & \enter{Model, e.g.\ Rigol DS1054Z} & \enter{S/N} \\
    Oscilloscope probe (x10)    & \enter{Model}                      & \enter{S/N} \\
    BNC cable (50\,\si{\ohm})   & \enter{Model}                      & ---         \\
    DC power supply             & \SI{10}{\volt} / \SI{5}{\volt}     & \enter{S/N} \\
    Experimentation board       & ---                                & ---         \\
    NE555 timer IC              & NE555                              & ---         \\
    CMOS NAND IC                & 74C00                              & ---         \\
    Quartz crystal              & \SI{12}{\mega\hertz}               & ---         \\
    Resistors                   & see individual tasks               & ---         \\
    Capacitors                  & see individual tasks               & ---         \\
    \bottomrule
  \end{tabular}
\end{table}

% ─────────────────────────────────────────────────────────────────────────────
\section{Theoretical Background}
% ─────────────────────────────────────────────────────────────────────────────

\subsection{NE555 Internal Structure}

The NE555 contains three equal resistors forming a voltage divider that sets the
reference levels at $\tfrac{1}{3}V_{CC}$ and $\tfrac{2}{3}V_{CC}$, two comparators,
an SR flip-flop, an npn discharge transistor (pin~7), and a push-pull output stage
(pin~3).

\subsection{Monostable Mode}

In monostable mode, the capacitor $C$ charges through $R_1$ after a falling trigger
edge at pin~2 pulls the output high. The output returns low once $V_C$ reaches
$\tfrac{2}{3}V_{CC}$, which gives the pulse duration:
\begin{equation}
  t_H = 1.1 \cdot R_1 \cdot C
  \label{eq:mono_theory}
\end{equation}

\subsection{Astable Mode}

In astable mode, $C$ charges through $R_1 + R_2$ and discharges through $R_2$ alone.
The resulting high and low times are:
\begin{align}
  t_H &= 0.693\,(R_1 + R_2)\,C \label{eq:astable_th_theory}\\
  t_L &= 0.693\,R_2\,C         \label{eq:astable_tl_theory}\\
  T   &= 0.693\,(R_1 + 2R_2)\,C \label{eq:astable_T_theory}
\end{align}

\subsection{RC Oscillator with CMOS Inverters}

Two cascaded inverters with a series resistor $R$ and capacitor $C$ in the feedback
path form a relaxation oscillator. The threshold voltages of the CMOS gates determine
the switching points, giving an oscillation frequency approximately:
\begin{equation}
  f \approx \frac{1}{2.2\,R\,C}
  \label{eq:inv_osc_theory}
\end{equation}

\subsection{Quartz Crystal Oscillator}

A quartz crystal behaves as a very high-Q series/parallel resonant circuit. When
connected in the feedback loop of an inverter amplifier, it forces oscillation at
its resonant frequency, providing frequency stability in the range of
$\pm10^{-5}$ to $\pm10^{-6}$.

% ─────────────────────────────────────────────────────────────────────────────
\section{Procedures}
% ─────────────────────────────────────────────────────────────────────────────

All experiments were carried out according to the laboratory instruction sheet
\textit{Lab\_01\_SwitchingCircuits}, Rev.\,1.1, dated 21/22.04.2026. Circuits were
assembled on the experimentation board. A supply voltage of \SI{10}{\volt} was used
for tasks~3.1--3.3, and \SI{5}{\volt} for task~3.4 (quartz oscillator).

Oscilloscope measurements were performed in two configurations where applicable:
\begin{enumerate}[leftmargin=1.5em]
  \item Using 10:1 oscilloscope probes (input impedance \SI{10}{\mega\ohm}
        $\|$ \SI{10}{\pico\farad}).
  \item Using direct 50\,\si{\ohm} BNC cables (task~3.1.4 only).
\end{enumerate}

Screenshots were saved for each relevant waveform. Component values were verified
with a multimeter before assembly where possible.

% ─────────────────────────────────────────────────────────────────────────────
\section{Results and Discussion}
% ─────────────────────────────────────────────────────────────────────────────

% ════════════════════════════════════════════════════════════════════
\subsection{Task 3.1 -- Monostable Multivibrator (NE555)}
% ════════════════════════════════════════════════════════════════════

\textbf{Component values used:}
$C = \SI{1}{\micro\farad}$,\quad
$R_1 = \SI{1}{\mega\ohm}$,\quad
$V_{CC} = \SI{10}{\volt}$.

% ── 3.1.1 ────────────────────────────────────────────────────────────
\subsubsection{Task 3.1.1 -- Calculation of Pulse Duration}

Applying equation~\eqref{eq:mono_theory}:
\begin{align*}
  t_{H,\text{calc}} &= 1.1 \times \SI{1}{\mega\ohm} \times \SI{1}{\micro\farad} \\
                    &= \SI{1.1}{\second}
\end{align*}
\[
  \boxed{t_{H,\text{calc}} = \SI{1.1}{\second}}
\]

% ── 3.1.2 ────────────────────────────────────────────────────────────
\subsubsection{Task 3.1.2 -- Measurement with Oscilloscope Probes}

Both the output signal $y$ (pin~3) and the capacitor voltage $V_C$ were measured
simultaneously using 10:1 oscilloscope probes.

\begin{table}[H]
  \centering
  \caption{Measured values -- monostable multivibrator, oscilloscope probes.}
  \label{tab:mono_probe}
  \begin{tabular}{lccc}
    \toprule
    Quantity & Symbol & Calculated & Measured \\
    \midrule
    Pulse duration              & $t_H$        & \SI{1.1}{\second}  & \val{\second} \\
    Peak capacitor voltage      & $V_{C,\max}$ & \SI{6.67}{\volt}   & \val{\volt} \\
    Rise time (10\,\%--90\,\%)  & $t_r$        & ---                & \val{\milli\second} \\
    \bottomrule
  \end{tabular}
\end{table}

\begin{figure}[H]
  \centering
  % Replace the placeholder box with your screenshot:
  % \includegraphics[width=0.85\linewidth]{fig/3_1_2_probes.png}
  \fbox{\rule{0pt}{5.5cm}\rule{12cm}{0pt}}
  \caption{Output $y$ (pin~3, upper trace) and capacitor voltage $V_C$ (lower trace)
           measured with 10:1 probes -- monostable multivibrator.}
  \label{fig:mono_probe}
\end{figure}

% ── 3.1.3 ────────────────────────────────────────────────────────────
\subsubsection{Task 3.1.3 -- Explanation of Differences}

\hint{
  \begin{itemize}[leftmargin=1.2em, itemsep=1pt]
    \item Resistor tolerance: $\pm5\,\%$; capacitor tolerance: $\pm20\,\%$
    \item NE555 internal resistance on pin~7 is not zero ($\approx$ a few \si{\ohm}--\SI{100}{\ohm})
    \item Comparator threshold not exactly $\tfrac{2}{3}V_{CC}$ (data sheet: $\pm15\,\%$)
    \item Leakage current of capacitor shifts the effective time constant
    \item Probe capacitance (\SI{10}{\pico\farad}) negligible for \SI{1}{\micro\farad}
  \end{itemize}
}

The measured pulse duration deviates from the calculated value by
$\Delta t_H = \val{\second}$, corresponding to a relative error of $\delta = \val{\percent}$.

\enter{Discuss which tolerance dominates (likely capacitor) and whether the deviation
is within the expected tolerance band.}

% ── 3.1.4 ────────────────────────────────────────────────────────────
\subsubsection{Task 3.1.4 -- Measurement with BNC Cables}

The measurement from task~3.1.2 was repeated, replacing the oscilloscope probes with
direct BNC cables (50\,\si{\ohm} input impedance).

\begin{table}[H]
  \centering
  \caption{Comparison: oscilloscope probes vs.\ BNC cables -- monostable multivibrator.}
  \label{tab:mono_bnc}
  \begin{tabular}{lccc}
    \toprule
    Quantity & Calculated & Probes & BNC cable \\
    \midrule
    Pulse duration $t_H$   & \SI{1.1}{\second} & \val{\second} & \val{\second} \\
    Peak $V_{C,\max}$      & \SI{6.67}{\volt}  & \val{\volt}   & \val{\volt}   \\
    \bottomrule
  \end{tabular}
\end{table}

\begin{figure}[H]
  \centering
  % \includegraphics[width=0.85\linewidth]{fig/3_1_4_bnc.png}
  \fbox{\rule{0pt}{5.5cm}\rule{12cm}{0pt}}
  \caption{Output $y$ and $V_C$ measured via BNC cable -- monostable multivibrator.}
  \label{fig:mono_bnc}
\end{figure}

\hint{
  \begin{itemize}[leftmargin=1.2em, itemsep=1pt]
    \item BNC cable: 50\,\si{\ohm} oscilloscope input, vs.\ \SI{10}{\mega\ohm} for probes
    \item The 50\,\si{\ohm} load in parallel with $R_1 = \SI{1}{\mega\ohm}$ is negligible
          at the output, but loading the capacitor node changes the effective discharge path
    \item Capacitor discharges faster through the low BNC impedance -- shorter $t_H$
    \item Amplitude of $V_C$ is attenuated by the voltage divider formed by $R_1$ and
          the 50\,\si{\ohm} load
  \end{itemize}
}

\enter{Describe and explain the observed differences. Calculate the effective parallel
resistance and the expected new time constant if relevant.}

% ── 3.1.5 ────────────────────────────────────────────────────────────
\subsubsection{Task 3.1.5 -- Output-Only Measurement (No Capacitor Probing)}

The output signal $y$ at pin~3 was measured alone, without a probe on the capacitor node.

\begin{table}[H]
  \centering
  \caption{Output pulse duration with and without simultaneous capacitor probing.}
  \label{tab:mono_outputonly}
  \begin{tabular}{lcc}
    \toprule
    Quantity & With $V_C$ probed & Without $V_C$ probed \\
    \midrule
    Pulse duration $t_H$ & \val{\second} & \val{\second} \\
    \bottomrule
  \end{tabular}
\end{table}

\begin{figure}[H]
  \centering
  % \includegraphics[width=0.85\linewidth]{fig/3_1_5_outputonly.png}
  \fbox{\rule{0pt}{5.5cm}\rule{12cm}{0pt}}
  \caption{Output $y$ (pin~3) only -- monostable multivibrator, no capacitor probing.}
  \label{fig:mono_outputonly}
\end{figure}

\hint{
  \begin{itemize}[leftmargin=1.2em, itemsep=1pt]
    \item Probe capacitance (\SI{10}{\pico\farad}) adds to $C$ at the capacitor node
    \item Removing the probe from the capacitor node slightly decreases total capacitance,
          resulting in a slightly shorter $t_H$
    \item Effect is very small: $\SI{10}{\pico\farad} \ll \SI{1}{\micro\farad}$
    \item Output stage (pin~3) is low impedance, so probing there has negligible effect
  \end{itemize}
}

\enter{Quantify any observed difference and state whether it is within measurement
uncertainty.}

% ════════════════════════════════════════════════════════════════════
\subsection{Task 3.2 -- Astable Multivibrator (NE555)}
% ════════════════════════════════════════════════════════════════════

\textbf{Component values used:}
$R_1 = \SI{10}{\kilo\ohm}$,\quad
$R_2 = \SI{10}{\kilo\ohm}$,\quad
$C = \SI{100}{\nano\farad}$,\quad
$V_{CC} = \SI{10}{\volt}$.

% ── 3.2.1 ────────────────────────────────────────────────────────────
\subsubsection{Task 3.2.1 -- Calculation of Oscillation Period}

Applying equations~\eqref{eq:astable_th_theory}--\eqref{eq:astable_T_theory}:
\begin{align*}
  t_H &= 0.693 \times (\SI{10}{\kilo\ohm} + \SI{10}{\kilo\ohm}) \times \SI{100}{\nano\farad}
       = \SI{1.386}{\milli\second} \\
  t_L &= 0.693 \times \SI{10}{\kilo\ohm} \times \SI{100}{\nano\farad}
       = \SI{0.693}{\milli\second} \\
  T   &= t_H + t_L = \SI{2.079}{\milli\second} \\
  f   &= \frac{1}{T} = \SI{481.0}{\hertz} \\
  D   &= \frac{t_H}{T} = \frac{R_1+R_2}{R_1+2R_2} = \SI{66.7}{\percent}
\end{align*}
\[
  \boxed{T_{\text{calc}} = \SI{2.079}{\milli\second}, \quad
         f_{\text{calc}} = \SI{481}{\hertz}, \quad
         D_{\text{calc}} = \SI{66.7}{\percent}}
\]

% ── 3.2.2 ────────────────────────────────────────────────────────────
\subsubsection{Task 3.2.2 -- Measurement}

\begin{table}[H]
  \centering
  \caption{Measured values -- astable multivibrator (NE555).}
  \label{tab:astable_ne555}
  \begin{tabular}{lccc}
    \toprule
    Quantity & Symbol & Calculated & Measured \\
    \midrule
    High time           & $t_H$        & \SI{1.386}{\milli\second} & \val{\milli\second} \\
    Low time            & $t_L$        & \SI{0.693}{\milli\second} & \val{\milli\second} \\
    Period              & $T$          & \SI{2.079}{\milli\second} & \val{\milli\second} \\
    Frequency           & $f$          & \SI{481}{\hertz}          & \val{\hertz}        \\
    Duty cycle          & $D$          & \SI{66.7}{\percent}       & \val{\percent}      \\
    Capacitor max.      & $V_{C,\max}$ & \SI{6.67}{\volt}          & \val{\volt}         \\
    Capacitor min.      & $V_{C,\min}$ & \SI{3.33}{\volt}          & \val{\volt}         \\
    \bottomrule
  \end{tabular}
\end{table}

\begin{figure}[H]
  \centering
  % \includegraphics[width=0.85\linewidth]{fig/3_2_2_astable.png}
  \fbox{\rule{0pt}{5.5cm}\rule{12cm}{0pt}}
  \caption{Output $y$ (pin~3, upper trace) and capacitor voltage $V_C$ (lower trace)
           -- astable multivibrator, NE555.}
  \label{fig:astable_ne555}
\end{figure}

% ── 3.2.3 ────────────────────────────────────────────────────────────
\subsubsection{Task 3.2.3 -- Explanation of Differences}

\hint{
  \begin{itemize}[leftmargin=1.2em, itemsep=1pt]
    \item Resistor tolerance $\pm5\,\%$, capacitor tolerance $\pm20\,\%$
    \item Internal NE555 discharge transistor $V_{CE,sat} > 0$ raises effective $R_2$
    \item Comparator thresholds deviate from ideal $\tfrac{1}{3}$ and $\tfrac{2}{3}$
    \item Probe adds parasitic capacitance to the $C$ node
    \item Since $R_1 = R_2$, the theoretical duty cycle is $\tfrac{2}{3}$;
          a diode in parallel with $R_2$ could yield 50\,\%
  \end{itemize}
}

\enter{State the absolute and relative deviations for $T$ and $f$. Identify the
dominant error source.}

% ── 3.2.4 ────────────────────────────────────────────────────────────
\subsubsection{Task 3.2.4 -- Rising Edge Zoom}

\begin{table}[H]
  \centering
  \caption{Rising edge parameters -- astable NE555 output.}
  \label{tab:astable_edge}
  \begin{tabular}{lcc}
    \toprule
    Quantity & Symbol & Measured \\
    \midrule
    Rise time (10\,\%--90\,\%) & $t_r$   & \val{\micro\second} \\
    Overshoot                  & ---     & \val{\percent}      \\
    \bottomrule
  \end{tabular}
\end{table}

\begin{figure}[H]
  \centering
  % \includegraphics[width=0.85\linewidth]{fig/3_2_4_edge.png}
  \fbox{\rule{0pt}{5cm}\rule{12cm}{0pt}}
  \caption{Zoomed rising edge of the NE555 astable output signal.}
  \label{fig:astable_edge}
\end{figure}

\hint{
  \begin{itemize}[leftmargin=1.2em, itemsep=1pt]
    \item NE555 output stage is a BJT totem-pole; switching from saturation to
          cut-off takes finite time (storage time)
    \item Capacitive load at pin~3 (scope probe, wiring) slows the edge
    \item Possible overshoot/ringing due to parasitic inductance of wires
    \item Rising edge typically faster than falling edge due to asymmetric BJT structure
  \end{itemize}
}

\enter{Describe the waveform shape. Is there ringing or overshoot? Relate the
observation to the internal BJT output stage and parasitic reactances.}

% ════════════════════════════════════════════════════════════════════
\subsection{Task 3.3 -- Astable Multivibrator Using Two Inverters (74C00)}
% ════════════════════════════════════════════════════════════════════

\textbf{Component values used:}
$R_S = \SI{1}{\kilo\ohm}$,\quad
$R = \SI{10}{\kilo\ohm}$,\quad
$C = \SI{1}{\nano\farad}$,\quad
$V_{CC} = \SI{10}{\volt}$.

\textbf{Expected frequency} from equation~\eqref{eq:inv_osc_theory}:
\[
  f_{\text{calc}} \approx \frac{1}{2.2 \times \SI{10}{\kilo\ohm} \times \SI{1}{\nano\farad}}
  = \SI{45.5}{\kilo\hertz}
\]

% ── 3.3.1 ────────────────────────────────────────────────────────────
\subsubsection{Task 1.1 -- Circuit Function}

\begin{figure}[H]
  \centering
  % \includegraphics[width=0.85\linewidth]{fig/3_3_1_function.png}
  \fbox{\rule{0pt}{5.5cm}\rule{12cm}{0pt}}
  \caption{Voltage at node P (lower trace) and output $U_{\text{OUT}}$ (upper trace)
           -- two-inverter oscillator.}
  \label{fig:inv_function}
\end{figure}

\hint{
  \begin{itemize}[leftmargin=1.2em, itemsep=1pt]
    \item First inverter output feeds back through $R$--$C$ to its own input at node~S
    \item When $U_{OUT}$ switches high, $C$ begins charging through $R$; voltage at S
          rises exponentially until it crosses the inverter switching threshold
    \item Second inverter buffers and re-inverts the signal
    \item $R_S$ limits current when $V_S$ overshoots $V_{CC}$ or undershoots GND
    \item Output is an approximate square wave; node P shows an exponential waveform
          clipped by the CMOS input protection diodes
  \end{itemize}
}

\enter{Describe what is observed on the oscilloscope at $U_{\text{OUT}}$ and P.
State the observed frequency.}

% ── 3.3.2 ────────────────────────────────────────────────────────────
\subsubsection{Task 1.2 -- Measurement at S and P}

\begin{table}[H]
  \centering
  \caption{Measured values -- two-inverter astable oscillator.}
  \label{tab:inv_osc}
  \begin{tabular}{lccc}
    \toprule
    Quantity & Symbol & Calculated & Measured \\
    \midrule
    Oscillation frequency & $f$          & \SI{45.5}{\kilo\hertz}   & \val{\kilo\hertz}   \\
    Period                & $T$          & \SI{22.0}{\micro\second} & \val{\micro\second} \\
    Voltage swing at S    & $\Delta V_S$ & $\approx V_{CC}$         & \val{\volt}         \\
    Voltage swing at P    & $\Delta V_P$ & $> V_{CC}$               & \val{\volt}         \\
    Peak voltage at P     & $V_{P,\max}$ & ---                      & \val{\volt}         \\
    Valley voltage at P   & $V_{P,\min}$ & ---                      & \val{\volt}         \\
    \bottomrule
  \end{tabular}
\end{table}

\begin{figure}[H]
  \centering
  % \includegraphics[width=0.85\linewidth]{fig/3_3_2_SP.png}
  \fbox{\rule{0pt}{5.5cm}\rule{12cm}{0pt}}
  \caption{Waveforms at S (upper trace) and P (lower trace) -- two-inverter oscillator.}
  \label{fig:inv_SP}
\end{figure}

\hint{
  \begin{itemize}[leftmargin=1.2em, itemsep=1pt]
    \item S: square wave switching between \SI{0}{\volt} and $V_{CC}$
    \item P: exponential charging/discharging, clamped by CMOS protection diodes
          to approximately $-0.7$\,V and $V_{CC}+0.7$\,V
    \item The CMOS switching threshold is typically at $\tfrac{1}{2}V_{CC}$
    \item The clipping of the P waveform is clearly visible on the oscilloscope
  \end{itemize}
}

\enter{Describe both waveforms in detail. Explain why the voltage at P exceeds the
supply rails and how the protection diodes limit it.}

% ── 3.3.3 ────────────────────────────────────────────────────────────
\subsubsection{Task 1.3 -- Purpose of $R_S$}

\hint{
  \begin{itemize}[leftmargin=1.2em, itemsep=1pt]
    \item Without $R_S$: unlimited current through the CMOS input protection diode
          when $V_P$ exceeds $V_{CC}$ or drops below GND -- IC damage possible
    \item $R_S$ limits diode current to a safe value:
          $I_{D} \approx \frac{V_{P,\max} - V_{CC}}{R_S}$
    \item $R_S$ also slightly reduces the oscillation frequency (adds to $R$)
    \item Typical recommendation: $R_S \geq \SI{1}{\kilo\ohm}$ for 74C-series CMOS
  \end{itemize}
}

With $R_S = \SI{1}{\kilo\ohm}$, the maximum protection diode current is limited to:
\[
  I_{D,\max} \approx \frac{V_{P,\max} - V_{CC}}{R_S}
  = \frac{\val{\volt} - \SI{10}{\volt}}{\SI{1}{\kilo\ohm}}
  = \val{\milli\ampere}
\]

\enter{Explain in your own words why $R_S$ is necessary and what would happen
without it.}

% ════════════════════════════════════════════════════════════════════
\subsection{Task 3.4 -- Quartz Crystal Oscillator (74C00, 12\,MHz)}
% ════════════════════════════════════════════════════════════════════

\textbf{Component values used:}
$R_1 = \SI{1}{\mega\ohm}$,\quad
$R_2 = \SI{220}{\ohm}$,\quad
$C_1 = C_2 = \SI{68}{\pico\farad}$,\quad
quartz: \SI{12}{\mega\hertz},\quad
$V_{CC} = \SI{5}{\volt}$.

% ── 3.4.1 ────────────────────────────────────────────────────────────
\subsubsection{Task 1.1 -- Output Voltage Measurement at $V_{CC} = \SI{5}{\volt}$}

\begin{table}[H]
  \centering
  \caption{Measured values -- quartz crystal oscillator at \SI{5}{\volt}.}
  \label{tab:quartz_5V}
  \begin{tabular}{lccc}
    \toprule
    Quantity & Symbol & Expected & Measured \\
    \midrule
    Frequency            & $f$      & \SI{12.000}{\mega\hertz} & \val{\mega\hertz}  \\
    Period               & $T$      & \SI{83.33}{\nano\second} & \val{\nano\second} \\
    Peak-to-peak voltage & $V_{pp}$ & $\leq V_{CC}$            & \val{\volt}        \\
    Rise time            & $t_r$    & ---                      & \val{\nano\second} \\
    \bottomrule
  \end{tabular}
\end{table}

\begin{figure}[H]
  \centering
  % \includegraphics[width=0.85\linewidth]{fig/3_4_1_quartz_5V.png}
  \fbox{\rule{0pt}{5cm}\rule{12cm}{0pt}}
  \caption{Output waveform of the quartz crystal oscillator at $V_{CC} = \SI{5}{\volt}$.}
  \label{fig:quartz_5V}
\end{figure}

% ── 3.4.2 ────────────────────────────────────────────────────────────
\subsubsection{Task 1.2 -- Explanation of the Curve}

\hint{
  \begin{itemize}[leftmargin=1.2em, itemsep=1pt]
    \item Quartz crystal: very high Q-factor ($Q \sim 10^4$--$10^6$), acts as a narrow
          bandpass filter in the feedback path
    \item $R_1$ biases the inverter into its linear (amplifying) region
    \item $C_1$, $C_2$ and the quartz form the Pierce oscillator topology
    \item Frequency determined almost entirely by the quartz -- not by $R$ or $C$ values
    \item $R_2$ damps parasitic oscillations and protects the output stage from
          excessive current
    \item Waveform near the crystal node may appear nearly sinusoidal due to filtering
  \end{itemize}
}

\enter{Describe the waveform shape (sinusoidal, square, or intermediate?). Comment on
the frequency accuracy compared to the quartz nominal frequency.}

% ── 3.4.3 ────────────────────────────────────────────────────────────
\subsubsection{Task 1.3 -- Effect of Increased Supply Voltage}

\begin{table}[H]
  \centering
  \caption{Output behaviour at different supply voltages -- quartz oscillator.}
  \label{tab:quartz_voltages}
  \begin{tabular}{lcccp{4.5cm}}
    \toprule
    $V_{CC}$ & $f$ (MHz) & $V_{pp}$ (V) & Waveform shape & Observations \\
    \midrule
    \SI{5}{\volt}  & \valx & \valx & \enter{shape} & Nominal operation \\
    \SI{10}{\volt} & \valx & \valx & \enter{shape} & \enter{describe} \\
    \SI{15}{\volt} & \valx & \valx & \enter{shape} & \enter{describe} \\
    \bottomrule
  \end{tabular}
\end{table}

\begin{figure}[H]
  \centering
  % \includegraphics[width=0.85\linewidth]{fig/3_4_3_voltages.png}
  \fbox{\rule{0pt}{5cm}\rule{12cm}{0pt}}
  \caption{Quartz oscillator output at \SI{5}{\volt}, \SI{10}{\volt}, and \SI{15}{\volt}.}
  \label{fig:quartz_voltages}
\end{figure}

\hint{
  \begin{itemize}[leftmargin=1.2em, itemsep=1pt]
    \item 74C00 absolute maximum $V_{CC}$: check datasheet (typically \SI{18}{\volt})
    \item Higher $V_{CC}$: faster propagation delay, hence slight frequency increase
    \item Higher $V_{CC}$: larger $V_{pp}$ at output
    \item Risk of IC damage above absolute maximum rating
    \item Quartz resonant frequency is nearly independent of $V_{CC}$, but very
          sensitive to temperature
  \end{itemize}
}

\enter{Describe and explain the effect on frequency, amplitude, and waveform shape.
Did the circuit still oscillate at all supply voltages tested?}

% ─────────────────────────────────────────────────────────────────────────────
\section{Conclusions}
% ─────────────────────────────────────────────────────────────────────────────

The following table provides a concise overview of calculated vs.\ measured values
for all tasks:

\begin{table}[H]
  \centering
  \caption{Summary: calculated vs.\ measured values for all tasks.}
  \label{tab:summary}
  \begin{tabular}{llccc}
    \toprule
    Task & Quantity & Calculated & Measured & Rel.\ error \\
    \midrule
    3.1 & $t_H$ (monostable)     & \SI{1.100}{\second}       & \val{\second}       & \val{\percent} \\
    3.2 & $T$ (astable NE555)    & \SI{2.079}{\milli\second} & \val{\milli\second} & \val{\percent} \\
    3.2 & $f$ (astable NE555)    & \SI{481}{\hertz}          & \val{\hertz}        & \val{\percent} \\
    3.3 & $f$ (inverter osc.)    & \SI{45.5}{\kilo\hertz}    & \val{\kilo\hertz}   & \val{\percent} \\
    3.4 & $f$ (quartz osc.)      & \SI{12.000}{\mega\hertz}  & \val{\mega\hertz}   & \val{\percent} \\
    \bottomrule
  \end{tabular}
\end{table}

\hint{
  \begin{itemize}[leftmargin=1.2em, itemsep=1pt]
    \item Compare all calculated vs.\ measured values from the table above
    \item Discuss the dominant error source for each circuit
    \item Comment on the probe vs.\ BNC measurement (task~3.1.4)
    \item State whether all experiment objectives were achieved
  \end{itemize}
}

\enter{Write 4--6 sentences summarising the main findings. Refer to Table~\ref{tab:summary}.}

% ─────────────────────────────────────────────────────────────────────────────
\section{Summary}
% ─────────────────────────────────────────────────────────────────────────────

\enter{3--5 sentence recap: restate the objective, briefly name the methods, state
the key numerical results, and give the main conclusion. No new information may appear
here. Do not copy from the Introduction -- paraphrase.}

% ─────────────────────────────────────────────────────────────────────────────
\section{Outlook}
% ─────────────────────────────────────────────────────────────────────────────

\enter{Suggest 2--3 possible follow-up experiments or improvements, for example:
duty-cycle adjustment of the astable NE555 using a diode in parallel with $R_2$;
temperature dependence of the RC-based oscillators;
comparison of quartz vs.\ ceramic resonator frequency stability.}

% ─────────────────────────────────────────────────────────────────────────────
\begin{thebibliography}{9}
% ─────────────────────────────────────────────────────────────────────────────

\bibitem{Floyd2006}
T.\,Floyd,
\textit{Digital Fundamentals}, 9th~ed.
Pearson Prentice Hall, 2006, ch.~7-6.

\bibitem{TietzeSchenk1999}
U.\,Tietze and C.\,Schenk,
\textit{Halbleiterschaltungstechnik}, 11.~Auflage.
Springer, 1999, Kap.~6.5.

\bibitem{ReportGuide}
C.\,Thorrold,
\textit{Guide for Writing Lab Reports}, Version~1.0.

\bibitem{NE555_DS}
Texas Instruments,
\textit{NE555 Precision Timers -- Datasheet},
\url{https://www.ti.com/lit/ds/symlink/ne555.pdf}, accessed April 2026.

\bibitem{74C00_DS}
Texas Instruments,
\textit{CD74C00 Quad 2-Input NAND Gate -- Datasheet},
\url{https://www.ti.com/lit/ds/symlink/cd74hc00.pdf}, accessed April 2026.

\end{thebibliography}

% ─────────────────────────────────────────────────────────────────────────────
\appendix
\section{Appendix A -- Raw Oscilloscope Screenshots}
% ─────────────────────────────────────────────────────────────────────────────

\enter{Insert full-resolution screenshots here, labelled with task number,
channel settings (V/div, time/div), and trigger settings.}

\section{Appendix B -- Actual Component Values}
% ─────────────────────────────────────────────────────────────────────────────

\begin{table}[H]
  \centering
  \caption{Component values measured with multimeter before assembly.}
  \label{tab:components}
  \begin{tabular}{lcccc}
    \toprule
    Component & Nominal & Measured & Tolerance & Task \\
    \midrule
    $R_1$ (monostable)  & \SI{1}{\mega\ohm}    & \val{\mega\ohm}    & $\pm5\,\%$  & 3.1 \\
    $C$  (monostable)   & \SI{1}{\micro\farad} & \val{\micro\farad} & $\pm20\,\%$ & 3.1 \\
    $R_1$ (astable)     & \SI{10}{\kilo\ohm}   & \val{\kilo\ohm}    & $\pm5\,\%$  & 3.2 \\
    $R_2$ (astable)     & \SI{10}{\kilo\ohm}   & \val{\kilo\ohm}    & $\pm5\,\%$  & 3.2 \\
    $C$  (astable)      & \SI{100}{\nano\farad}& \val{\nano\farad}  & $\pm20\,\%$ & 3.2 \\
    $R_S$ (inverter)    & \SI{1}{\kilo\ohm}    & \val{\kilo\ohm}    & $\pm5\,\%$  & 3.3 \\
    $R$  (inverter)     & \SI{10}{\kilo\ohm}   & \val{\kilo\ohm}    & $\pm5\,\%$  & 3.3 \\
    $C$  (inverter)     & \SI{1}{\nano\farad}  & \val{\nano\farad}  & $\pm20\,\%$ & 3.3 \\
    \bottomrule
  \end{tabular}
\end{table}

\end{document}